\begin{abstract}
A receiver-derived current space is a useful starting point for electromagnetic inverse scattering, but its modes are ranked by reception rather than by the material update they produce. We identify the material consequence of a specified current-space modification in vector Maxwell equations. Eliminating retained currents exposes three factors: material injection, multiple-scattering feedback, and receiver return. Their two material pullbacks contain the change in the regularized Gauss--Newton (GN) step. Its usefulness depends jointly on curvature change and a residual-dependent offset; increased information magnitude need not improve the step. We therefore start from the receiver-derived basis and test conditional one-out/one-in exchanges at fixed current dimension. A richer reduced anchor ranks proposals, and full directional tangent actions verify at most two finalists per round against the unchanged full quadratic. Acceptance uses an online numerical significance tolerance, with no full optimal step or truth access. Original four-object experiments reduce the median object-summed full-GN gap by 65.2\%. With the rule locked, all ten complete additional objects improve, with a median reduction of 50.7\%; one further object retains a late-state prerequisite failure. Directly dark currents, information-positive but utility-negative replacements, and pair interactions demonstrate why reception and information summaries alone are insufficient. The resulting contribution is a Maxwell mechanism and a bounded conditional design for material-step fidelity, with explicit acquisition costs and separate noisy and nonlinear transfer tests.
\end{abstract}
